\begin{table}[!htbp]
\caption{Practical decision framework for evaluating an added urban data
source in an integer-count forecasting system.}
\label{tab:practical_decision_framework}
\centering
\setlength{\tabcolsep}{3.5pt}
\begin{tabularx}{\textwidth}{
>{\RaggedRight\arraybackslash}p{2.75cm}
>{\RaggedRight\arraybackslash}X
>{\RaggedRight\arraybackslash}X
>{\RaggedRight\arraybackslash}p{3.15cm}}
\toprule
\textbf{Decision question} & \textbf{Required evidence} &
\textbf{Operational response} & \textbf{Boundary}\\
\midrule
Does the source create an estimated prediction opportunity? &
DELB and CMI point estimates in frozen target populations &
Retain the source for calibrated evaluation if the estimated floor
declines materially & Point estimates alone are insufficient for
source attribution or realized benefit\\
Is the apparent gain distinguishable from finite-sample structure? &
Support diagnostics, known-truth calibration, and design-matched
randomization & Attribute the gain only to the tested design and only when
calibration and power support separation & Low-powered non-rejection is
inconclusive; randomization provides no causal identification\\
Can available models use the source on unseen data? &
Matched held-out models with identical targets, samples, tuning rules, and
loss & Deploy only city--model combinations showing stable held-out benefit;
monitor the Predictability Gap & Model-specific gains require separate
validation across cities and model families\\
Does value persist outside the service footprint? &
Training-period-frozen coverage-conditioned and complete crime-supported
deployment populations &
Report both populations and treat coverage as part of the data product &
Service-area estimates apply only to covered locations\\
\bottomrule
\end{tabularx}
\end{table}
